\documentclass[../main.tex]{subfiles}
\begin{document}
\section{Reduction to a $2$-adic Question}
For each $n \in \Z_{> 1}$ denote 
\begin{align*}
  & f_n\left( x \right) = x^{n-1} - \sum_{i=0}^{n-2} x^i \\
  & g_n\left( x \right) = \left( x - 1 \right)f_n\left( x \right) = x^n - 2x^{n-1} + 1
\end{align*}
We will ommit for now the proof [or even any reference to that fact] that $f_n$ is irreducible over $\Q$.
Instead, we will reduce the question of whether its Galois group is the full symmetric group to a $2$-adic question.
Denote by $L_n$ the splitting field of $f_n$ over $\Q$ and by $G_n = \Gal\left( L_n / \Q \right)$ its Galois group.

\begin{lemma}
  Let $\P \subseteq \O_{L_n}$ be a prime at which $f_n$ ramifies that lies over an odd prime $\left( p \right) \subseteq \Z$.
  Then the inertia group of $f_n$ at $\P$ is generated by a single transposition.
\end{lemma}
\begin{proof}
  It suffices to show that over $\F_p$, the polynomial $f_n$ has exactly one double root and all other roots are simple.
  Indeed, assume that this is the case and denote by $\alpha$ the unique double root of $f_n$ over $\F_p$.
  Then, any other root $\beta$ of $f_n$ over $\F_p$ lifts to a single root, and hence any element of the inertia group must fix it.
  Therefore, the inertia group is embedded into the permutation group of the two lifts of $\alpha$, and hence is either trivial or generated by a single transposition.
  However, since by assumption $f_n$ ramifies at $p$, the inertia group is non-trivial, and hence it is indeed generated by the transposition of the two lifts of $\alpha$.
  
  We shall show then that $f_n$ has exactly one double root over $\F_p$ and all other roots are simple.
  If $\alpha$ is a multiple root of $f_n$ over $\F_p$, it is in particular a multiple root of $g_n$ and hence
  \[
    0 = g_n'\left( \alpha \right) = n \alpha^{n-1} - 2\left( n - 1 \right) \alpha^{n-2} 
  \]
  and hence, since $\alpha \neq 0$,
  \[
    \alpha = \frac{2\left( n - 1 \right)}{n}
  \]
  and it is in particular unique.
  
  Since $f_n$ is by assumption ramified at $p$, such $\alpha$ exists.
  We will show that the multiplicity of $\alpha$ is not greater than $2$.
  It suffices to show that $g_n$ has no root of multiplicity greater than $2$, and since all of its roots are non-zero, its multiset of roots is obtained from the multiset of $h_n\left( x \right) = x^n - 2x + 1$ by inversion, and hence it suffices to show that $h_n$ has no root of multiplicity greater than $2$ over $\F_p$.
  However, any root of multiplicity greater than $2$ must also be a root of $h_n''\left( x \right) = n\left( n - 1 \right)x^{n-2}$, and hence must be $0$, which is not a root of $h_n$.
\end{proof}
\begin{lemma}
  Let $L/ \Q$ be a Galois extension.
  Then 
  \[
    G = \sum_{\P \subseteq \O_L} I_{\P}
  \]
  where $I_{\P}$ is the inertia group at the prime $\P$.
\end{lemma}
\begin{proof}
  The only unramified extension of $\Q$ is $\Q$ itself, and hence the fixed field of $\sum_{\P} I_{\P}$ is $\Q$.
\end{proof}
\begin{lemma}
  A transitive permutation group generated by transpositions is the full symmetric group.
\end{lemma}
\begin{proof}
  Sketch inspired by Itamar Israeli:
  There is a seuqnece of transpositions that connects any two elements.
  Given a product of a pair of transpositions 
  \[
    \left( a \; b \right) \left( b \; c \right) = \left( a \; b \; c \right)
  \]
  it can be shoretned to 
  \[
    \left( b \; c \right) \left( a \; b \right) \left( b \; c \right) = \left( a \; c \right) 
  \]
\end{proof}
\begin{corollary}
  If $n$ is odd, then $G_n = S_{n-1}$.
\end{corollary}
\begin{corollary}
  Assume that $n$ is even.
  To prove that $G_n = S_{n-1}$, it suffices to show that for any prime $\P \subseteq \O_{L_n}$ above $2$, the inertia group $I_{\P}$ is generated by a transposition.
\end{corollary}

We can also derive a slightly different reduction. 
\begin{lemma}
  A primitive transitive permutation group $G$ containing a transposition is the full symmetric group. 
\end{lemma}
\begin{proof}
  Since it is transitive and primitive, any non-trivial normal subgroup of $G$ is also transitive.
  Indeed the orbits of any normal subgroup form a partition of the set preserved by $G$.
  Therefore, the normal subgroup generated by the transposition whose existence is assumed is transitive, and hence must be the full symmetric group, and it is generated by all the conjugates of that transposition, which are all transpositions in themselves. 
\end{proof}
\begin{lemma}
  If $n > 2$ then $d_n$ is not a power of $2$. 
\end{lemma}
\begin{proof}
  By \cite{MARTIN2004213},
  \[
    d_n = \pm 
    \frac{n^n - 2^n \left( n - 1 \right)^{ n - 1 } }{
      \left( n - 2 \right)^2
    }
  \]
  If $n$ is odd, then $n^n - 2^n \left( n - 1 \right)^{ n - 1 }$ is odd as well, and hence $d_n$ is not a power of $2$.
  Assume that $n$ is even, then by Lemma \ref{lem:2-adic-valuation-of-discriminant},
  \[
    v_2\left( d_n \right) = n - 2
  \] 
\end{proof}
\begin{corollary}
  To prove that $G_n = S_{n-1}$, it suffices to show that $G_n$ is primitive and that the discriminant of $L_d$ over $\Q$ is not a power of $2$ (for which is sufficient to show that $d_n$ is not a power of $2$).
\end{corollary}
\begin{lemma}
  A permutation group $G$ that contains an $n$-cycle and a transposition is symmetric.
\end{lemma}
\begin{proof}
  \editorial{I guess we could use a version of Itamar's proof here as well}
\end{proof}



\end{document}